\documentclass[11pt]{article}
\usepackage[T1]{fontenc}
\usepackage[utf8]{inputenc}
\usepackage{lmodern}
\usepackage[a4paper,margin=28mm]{geometry}
\usepackage{amsmath,amssymb,amsthm,mathtools,booktabs,array}
\usepackage[hidelinks]{hyperref}
\usepackage{microtype}
\newtheorem{theorem}{Theorem}[section]
\newtheorem{lemma}[theorem]{Lemma}
\newtheorem{proposition}[theorem]{Proposition}
\newtheorem{corollary}[theorem]{Corollary}
\theoremstyle{remark}
\newtheorem{remark}[theorem]{Remark}
\newcommand{\D}{\mathcal D}
\newcommand{\T}{\mathbb T}
\newcommand{\Z}{\mathbb Z}
\newcommand{\N}{\mathbb N}
\newcommand{\MC}{M_{\mathrm C}}
\newcommand{\MH}{M_{\mathrm H}}
\newcommand{\ee}[1]{\mathrm e(#1)}
\DeclareMathOperator{\diam}{diam}

\newcommand{\Obloc}{\mathcal O_3}
\newcommand{\Obreal}{\mathcal O_\infty}
\newcommand{\F}{\mathcal F}
\newcommand{\I}{\mathcal I}
\newcommand{\Sc}{\mathcal S}
\newcommand{\bits}{\{0,1\}}
\title{Integer quotients of ternary Cantor numbers:\\
a power lower bound and an infinite exceptional family}
\author{Timofei Iuzhakov}
\date{8 November 2026\\Version 2.0.0}
\begin{document}
\maketitle

\begin{abstract}
Let $\D$ be the nonnegative integers whose ternary digits belong to
$\{0,1\}$. We prove two results about positive integer quotients of
members of $\D$. First, the number $\MC(X)$ of integers $1\le M\le X$
with $N,MN\in\D$ for some positive integer $N$ satisfies
\[
 \MC(X)\gg X^{\beta_*},\qquad
 \beta_*={(1+s)^2\over7-5s}=0.691726633148796\ldots,
 \quad s=\log_3 2.
\]
Version 2 improves the original exponent $0.6660585297\ldots$,
without an epsilon loss, by an endpoint mean-discrepancy estimate
and an exact fixed-residue-class energy certificate.
The same bound holds for the number of positive-dimensional
intersections of the $3$-adic Cantor set with its inverse integer
dilates. Second, for every $k\ge0$, the integer
$621\cdot3^{4k}-20$ is not such a quotient, although it avoids the
standard lowest-digit and leading-digit obstructions. This proves
that the sporadic exception set of Bai, Meleshko, Riasat and Shallit
is infinite. The first proof combines Fourier sampling over
Cantor-restricted denominators with a finite-word Fibonacci estimate
and a sum over $3$-adic distances. The second constructs an explicit
predecessor-closed carry certificate of $38k+50$ states for $k\ge1$,
with a separate certificate for $k=0$. The original proofs and the new improvement are given in full.
\end{abstract}

\section{The quotient problem and the two main results}

\paragraph{Version history.}
Version 1 (6 November 2026) established Theorems~\ref{thm:power}
and~\ref{thm:family}. Version 2 (8 November 2026) retains their full
proofs and adds Theorem~\ref{ep:main} in Section~\ref{ep:section},
raising the exponent from $0.6660585297\ldots$ to
$0.691726633148796\ldots$ without an epsilon loss.
The article title, sole author, and citation key
\texttt{iuzhakov2026cantorquotients} are unchanged.
The sharp growth of $\MC$ and Erd\H{o}s Problem 406 remain open here.


Put
\[
 \D=\left\{\sum_{j\ge0}\epsilon_j3^j:
   \epsilon_j\in\{0,1\},\ \epsilon_j=0\text{ for all large }j\right\},
 \qquad \D^+=\D\setminus\{0\},
\]
and let
\[
 \Sigma=\left\{\sum_{j\ge0}\epsilon_j3^j:
                    \epsilon_j\in\{0,1\}\right\}\subset\Z_3.
\]
Here $\N$ denotes the positive integers. Define the set of positive
integer quotients and its counting function by
\[
 \mathcal V=\N\cap(\D^+/\D^+),\qquad
 \MC(X)=\#(\mathcal V\cap[1,X])\quad(X\ge1).
\]
Equivalently, $M\in\mathcal V$ if some positive integer $N$ has
$N,MN\in\D$. Positivity of the witness is essential: allowing $N=0$
would make this condition vacuous. For a positive integer $M$, also put
\[
 C(1,M)=\{x\in\Z_3:x,Mx\in\Sigma\},\qquad
 \MH(X)=\#\{1\le M\le X:\dim_H C(1,M)>0\}.
\]

\subsection{A power lower bound}

Lagarias asks for upper and lower bounds for these two counting
functions in Problems 1.10 and 1.11 of~\cite{Lagarias}, numbered
Problems 1 and 2 in the arXiv version. Taking $N=1$ gives the immediate
lower bound of order $X^{\log_3 2}$. A positive integer witness also
implies positive Hausdorff dimension by the nonoverlapping-block
construction in~\cite{Lagarias}; we include this construction below.
The inclusion of the integer-witness multipliers in the
positive-dimensional multipliers is strict in general.

\begin{theorem}\label{thm:power}
Let
\[
 s=\log_3 2,\qquad \varphi=\frac{1+\sqrt5}{2},\qquad
 g=\log_3\varphi,\qquad
 \alpha=\frac{s-g}{1+g},\qquad
 \beta=s+(2s-1)\alpha.
\]
There are absolute constants $c>0$ and $X_0$ such that
\[
 \MC(X)\ge cX^\beta,\qquad \MH(X)\ge cX^\beta
 \quad (X\ge X_0).
\]
In particular,
\[
 \alpha=0.1341512347\ldots,
 \qquad \beta=0.6660585297\ldots>\log_3 2.
\]
\end{theorem}

The result improves the exponent in the immediate lower bound. It does
not determine the sharp asymptotic size or positive density of either
counting function. The distribution estimates used in its proof are
proved directly, with their dependence on growing denominators explicit.

\subsection{An infinite exceptional family}

Following Bai, Meleshko, Riasat and Shallit~\cite[\S1.4]{BMRS}, define
\begin{align*}
 \Obloc&=\{N\in\N:N\equiv2\cdot3^{i-1}\pmod{3^i}
                             \text{ for some }i\ge1\},\\
 \Obreal&=\N\cap\bigcup_{i\ge0}
                    [\tfrac32\,3^i,\,2\cdot3^i],\\
 \F&=\N\setminus(\mathcal V\cup\Obloc\cup\Obreal).
\end{align*}
The first two sets encode the lowest nonzero digit and leading-digit
obstructions to being a quotient. They are denoted $\mathcal D$ and
$\mathcal E$ in~\cite{BMRS}; our different symbols avoid conflict with
$\D$, the digit-integer set. The original Cantor integers in that paper
have digits $0,2$. Their positive members form $2\D^+$, so their
positive quotient set is exactly $\D^+/\D^+$. The sporadic set $\F$
therefore agrees with theirs.

Bai et al. ask whether $\F$ is infinite and propose the family in the
following theorem.

\begin{theorem}\label{thm:family}
For every integer $k\ge0$,
\[
                  621\cdot3^{4k}-20\in\F.
\]
In particular, there are no $a,b\in\D^+$ with
$b=(621\cdot3^{4k}-20)a$, and $\F$ is infinite.
\end{theorem}

The case $k=4$ is verified in Jandhyala's versioned
preprint~\cite[\S6]{Jandhyala}, which leaves the all-$k$ assertion open.
Its main theorems concern the different quotient problem for zero-free
balanced digits $\{-1,1\}$. Our proof of Theorem~\ref{thm:family}
uses the identity $b+20a=23\cdot3^{4k+3}a$ to obtain a binary delay
register with a carry in a fixed set of $22$ integers. An explicit
predecessor-closed set separates the initial state from the only
nontrivial predecessor of an accepting state.

\subsection{Relation between the two arguments}

The quotient problem is also studied directly by Anders, Dawsey,
Reznick and Sisneros-Thiry~\cite{ADRS}. Multiplicative intersections
and their automata are treated in~\cite{AL,ABL}. The analytic part
of our counting proof is related to work on the distribution of
integers with missing digits~\cite{Konyagin,MS}, but no external
uniformity estimate is used without proof.

The two main theorems are logically independent. The first proves
abundance among all positive integers; the second gives specified
integers which fail to be quotients despite the elementary necessary
conditions. The exceptional family is not a matching upper bound for
$\MC(X)$, and its nonrepresentation does not by itself determine
$\dim_H C(1,M)$ for those multipliers.

Neither theorem concerns the number of exponents $n$ for which
$2^n\in\D$. In particular, neither proves the pointwise assertion in
Erd\H{o}s problem 406, excludes $1$ from the infinite-visit exceptional
set, or determines a finite-visit dimension limit. The arguments below
use neither a general causal-container theorem nor a finite-rank
unit-equation theorem. We make no claim of exhaustive priority
verification.

\section{Fourier sampling over reduced fractions}

For $l\ge0$, let
\[
 \D_l=\left\{\sum_{j=0}^{l-1}\epsilon_j3^j:
                     \epsilon_j\in\{0,1\}\right\},
 \quad
 F_l(t)=2^{-l}\sum_{z\in\D_l}\ee{zt},\quad f_l(t)=|F_l(t)|,
\]
where $\ee{t}=\exp(2\pi it)$ and $t\in\T=\mathbb R/\mathbb Z$.
The product identity gives
\begin{equation}\label{count:eq:product}
 F_l(t)=\prod_{j=0}^{l-1}\frac{1+\ee{3^jt}}2,
 \qquad
 f_{l+l'}(t)=f_l(t)f_{l'}(3^lt),\qquad 0\le f_l(t)\le1.
\end{equation}

\begin{lemma}[Elementary sampling]\label{count:lem:sampling}
Let $0<\delta\le1$ and let $T\subset\T$ be finite, with circular
separation at least $\delta$.
If $P(t)=\sum_{j=0}^K c_j\ee{jt}$, then
\[
 \sum_{t\in T}|P(t)|^2
 \le (\delta^{-1}+4\pi K)\sum_{j=0}^K|c_j|^2.
\]
\end{lemma}
\begin{proof}
Center pairwise disjoint half-open arcs $I_t$ of length $\delta$ at the
points of $T$. For $G=|P|^2$ and $u\in I_t$,
\[
 G(t)\le G(u)+\int_{I_t}|G'(v)|\,dv.
\]
Integrating in $u$ and summing in $t$ bounds the left side by
$\delta^{-1}\|P\|_2^2+\int_\T|G'|$.
Since $\int|G'|\le2\|P\|_2\|P'\|_2\le4\pi K\|P\|_2^2$,
Parseval proves the assertion.
\end{proof}

\begin{lemma}\label{count:lem:moments}
For every $h\ge0$,
\[
 \int_\T|F_h|^2=2^{-h},\qquad
 \int_\T|F_h|^4=(3/8)^h.
\]
The degrees of $F_h$ and $F_h^2$ are $(3^h-1)/2$ and $3^h-1$.
\end{lemma}
\begin{proof}
The second moment follows from the distinct frequencies in $\D_h$.
For the fourth moment, the coefficients of $F_h^2$ count sums of two
elements of $\D_h$. At each ternary position the sums $0,1,2$ have
multiplicities $1,2,1$. There is no carry, so the sum of the squares
of the unnormalized coefficients is $(1^2+2^2+1^2)^h=6^h$.
The normalization is $16^h$. The degree statements follow from
$\max\D_h=(3^h-1)/2$.
\end{proof}

For real $D\ge2$, write
\[
 T_D=\{v/d\bmod1:2\le d\le D,\ (d,3)=1,
                       \ 1\le v<d,\ (v,d)=1\}.
\]
Distinct points of $T_D$ have circular separation at least $D^{-2}$.

\begin{lemma}[Farey Fourier bound]\label{count:lem:farey}
Put $\eta=\log_3(27/16)<s$. There is an absolute constant $C$ such that,
if $D\ge2$ and $L\ge3\lceil2\log_3D\rceil$, then
\begin{equation}\label{count:eq:farey}
 \sum_{t\in T_D}f_L(t)\le CD^\eta.
\end{equation}
For each fixed $D$, the sum on the left tends to zero as $L\to\infty$.
\end{lemma}
\begin{proof}
Take $h=\lceil2\log_3D\rceil$, so $D^2\le3^h<3D^2$.
Lemmas \ref{count:lem:sampling} and \ref{count:lem:moments} give
\[
 \sum_{t\in T_D}f_h(t)^2\le C(3/2)^h,
 \qquad
 \sum_{t\in T_D}f_h(t)^4\le C(9/8)^h.
\]
Cauchy--Schwarz yields
\begin{equation}\label{count:eq:third}
 \sum_{t\in T_D}f_h(t)^3
 \le C(27/16)^{h/2}\le C'D^\eta.
\end{equation}
Let $p=\lfloor L/h\rfloor\ge3$. By \eqref{count:eq:product}, after discarding
the remainder of length $L-ph$, H\"older's inequality gives
\begin{align*}
 \sum_{t\in T_D}f_L(t)
 &\le\prod_{j=0}^{p-1}
   \left(\sum_{t\in T_D}f_h(3^{jh}t)^p\right)^{1/p}\\
 &=\sum_{t\in T_D}f_h(t)^p
 \le\sum_{t\in T_D}f_h(t)^3.
\end{align*}
Multiplication by $3^{jh}$ permutes $T_D$, since its reduced
denominators are coprime to $3$. This proves \eqref{count:eq:farey}.
For a fixed nonzero $v/d\in T_D$, the orbit under multiplication by $3$
is periodic and avoids zero. The product in \eqref{count:eq:product} therefore
tends to zero in modulus. Summing over the fixed finite set proves the
last assertion.
\end{proof}

\section{Distribution for most Cantor denominators}

For $Y\ge3$, set
\[
 S_Y=\{N\in\D:1\le N\le Y,\ N\equiv1\pmod3\}.
\]
Elementary digit counting gives
\begin{equation}\label{count:eq:SY}
 \tfrac14Y^s\le |S_Y|\le2Y^s.
\end{equation}
For the lower bound use all words of length $\lfloor\log_3Y\rfloor$
whose low digit is $1$. For the upper bound use all words of length
$K=\lfloor\log_3Y\rfloor+1$.

\begin{lemma}[Divisor multiplicity]\label{count:lem:divisors}
For integers $1\le d\le Y$,
\[
 T_Y(d):=\#\{N\in S_Y:d\mid N\}\le4(Y/d)^s.
\]
\end{lemma}
\begin{proof}
Put $j=\lfloor\log_3d\rfloor$. Fixing the digits at positions
$j,\ldots,K-1$ covers $\D_K$ by $2^{K-j}$ intervals, each of length
$(3^j-1)/2<d$. Each interval contains at most one multiple of $d$.
The inequalities $3^K\le3Y$ and $3^j>d/3$ give
$2^{K-j}<4(Y/d)^s$.
\end{proof}

For $3\nmid N$, define
\[
 R_N(L)=\sum_{v=1}^{N-1}f_L(v/N),\qquad R_1(L)=0.
\]

\begin{proposition}[Average discrepancy on restricted denominators]
\label{count:prop:mixing}
As $Y,L\to\infty$ subject to $L\ge3\lceil2\log_3Y\rceil$,
\begin{equation}\label{count:eq:mixing}
 \sum_{N\in S_Y}R_N(L)=o(Y^s).
\end{equation}
Consequently, for all sufficiently large such $Y,L$, at least
$Y^s/8$ elements $N\in S_Y$ satisfy $R_N(L)\le1/2$.
\end{proposition}
\begin{proof}
Reducing the fractions $v/N$ gives the exact identity
\[
 \sum_{N\in S_Y}R_N(L)
 =\sum_{\substack{2\le d\le Y\\(d,3)=1}}
   T_Y(d)\sum_{\substack{1\le v<d\\(v,d)=1}}f_L(v/d).
\]
Use dyadic bins $(2^{j-1},2^j]$ for $1\le j\le J=\lfloor\log_2Y\rfloor$,
followed by the possibly empty bin $(2^J,Y]$. By Lemmas
\ref{count:lem:divisors} and \ref{count:lem:farey}, the contribution of a complete
bin, divided by $Y^s$, is bounded by
\[
 C2^{-js}\sum_{t\in T_{2^j}}f_L(t)
 \le C'2^{-j(s-\eta)}.
\]
The latter bounds are summable, since $\eta<s$. For each fixed $j$,
the actual contribution tends to zero, by the last assertion of
Lemma \ref{count:lem:farey}. First make the tail of the summable bounds small,
and then take $L$ large enough for the finitely many remaining bins.
The final partial bin is at most $CY^{\eta-s}$ after division by $Y^s$,
and also tends to zero. This proves \eqref{count:eq:mixing}.
At most $2\sum_NR_N(L)=o(Y^s)$ denominators have $R_N(L)>1/2$;
combine this with \eqref{count:eq:SY}.
\end{proof}

\begin{proposition}[First moment]\label{count:prop:first}
Let $L=\lfloor\log_3X\rfloor$, and suppose $Y\to\infty$ with
$L\ge3\lceil2\log_3Y\rceil$ and $X^sY^{s-1}\to\infty$.
There is a subset $G\subset S_Y$ such that
\[
 \sum_{N\in G}|A_N(X)|\ge\tfrac1{64}X^sY^{2s-1},
 \qquad A_N(X)=\{1\le M\le X:NM\in\D\},
\]
for all sufficiently large $X$.
\end{proposition}
\begin{proof}
Choose $G=\{N\in S_Y:R_N(L)\le1/2\}$, which has size at least $Y^s/8$.
For each $N\in G$ put $L_N=\lfloor\log_3(NX)\rfloor\ge L$.
By \eqref{count:eq:product}, $R_N(L_N)\le R_N(L)\le1/2$.
Orthogonality of additive characters modulo $N$ gives
\[
 \#\{z\in\D_{L_N}:N\mid z\}
 \ge\frac{2^{L_N}}N(1-R_N(L_N))\ge\frac{2^{L_N-1}}N.
\]
Every such $z$ is at most $NX$. Removing zero and dividing by $N$ yields
\[
 |A_N(X)|\ge\tfrac14X^sN^{s-1}-1.
\]
Since $s<1$, summing and using \eqref{count:eq:SY} gives
\[
 \sum_{N\in G}|A_N(X)|
 \ge\tfrac1{32}X^sY^{2s-1}-2Y^s
 \ge\tfrac1{64}X^sY^{2s-1}.
\]
\end{proof}

\section{A finite Fibonacci bound for overlaps}

For $r\in1+3\Z_3$, $r\ne1$, put $e=v_3(r-1)\ge1$. Define
\[
 W_n(r)=\{x\in\D_n:rx\bmod3^n\in\D_n\},
\]
using representatives in $[0,3^n)$ for reduction. Let $H_n(e)$ be the
number of independent sets of the graph on $\{0,\ldots,n-1\}$ whose
edges join positions at distance $e$.

\begin{lemma}[Finite-word bound]\label{count:lem:golden}
For every $n\ge0$,
\[
 |W_n(r)|\le H_n(e)\le\varphi^{n+e}.
\]
The bound applies to all locally admissible words, including those
which have no infinite admissible continuation.
\end{lemma}
\begin{proof}
Write $r=1+3^eu$, where $u\in\Z_3^\times$. Write the input digits as
$x_j\in\{0,1\}$ and the digits of $ux$ as $y_j\in\{0,1,2\}$.
In the addition $x+3^eux$, let $\gamma_q\in\{0,1\}$ be the incoming
carry at position $q$. For $q\ge e$, put $B_q=y_{q-e}+\gamma_q$.
The allowed values and next carries are exactly
\[
\begin{array}{c|c|c}
 B_q&\text{allowed }x_q&\gamma_{q+1}\\\hline
 0&0,1&0\\1&0&0\\2&1&1\\3&0,1&1
\end{array}
\]
Thus $\gamma_{q+1}=\mathbf1_{B_q\ge2}$ on admissible paths,
independently of the chosen admissible value of $x_q$. For $q<e$ the
carry is zero.

Fix $x_{<i}$. The preceding identity determines $\gamma_{i+e}$ from
$y_0,\ldots,y_{i-1}$, hence from $x_{<i}$, independently of the
intermediate admissible digits. Also
$y_i\equiv u_0x_i+R_i(x_{<i})\pmod3$, with $u_0\in\{1,2\}$.
The output at position $i+e$ consequently has the form
\begin{equation}\label{count:eq:forcing}
 x_{i+e}+u_0x_i+Q_i(x_{<i})\pmod3.
\end{equation}
At least one of $x_i=0,1$ makes the offset in \eqref{count:eq:forcing}
nonzero modulo $3$, leaving at most one allowed value of $x_{i+e}$,
regardless of the intermediate admissible digits.

Mark that future position as forced. A mark means that each subsequent
admissible prefix on the branch has at most one child at the marked
position; it does not specify a digit common to different branches.
Maintain the set $S$ of future marks. At a marked position there is
at most one child. At an unmarked position with one child, retaining
$S$ suffices for an upper bound. At an unmarked position $i$ with two
children, one child retains $S$, while the other has the additional
mark $i+e$, if that position is less than $n$. Additional restrictions
may be discarded.

Backward induction now compares the number of completions with
independent sets of the remaining distance-$e$ graph, with positions
in $S$ required absent. At an unmarked vertex $i$, its deletion
recursion is the sum of the counts without and with the additional
mark $i+e$; at a marked vertex there is a single continuation.
Requiring additional positions absent can only decrease the count.
This proves $|W_n(r)|\le H_n(e)$, starting with no marks.

The graph is a union of at most $e$ paths, one in each residue class
modulo $e$. A path of $m$ vertices has $F_{m+2}$ independent sets,
and $F_{m+2}\le\varphi^{m+1}$ by the Fibonacci recurrence.
Multiplying these inequalities gives $H_n(e)\le\varphi^{n+e}$.
\end{proof}

\begin{lemma}[Valuation summation]\label{count:lem:valuation}
For $Y\ge3$,
\[
 \sum_{\substack{N<N'\\N,N'\in S_Y}}
       \varphi^{v_3(N'-N)}\ll Y^{2s}.
\]
\end{lemma}
\begin{proof}
Put $K=\lfloor\log_3Y\rfloor+1$. For fixed $N\in S_Y$ and $1\le e<K$,
\[
 \#\{N'\in S_Y:N'\equiv N\pmod{3^e}\}
 \le2^{K-e}\le2Y^s2^{-e}.
\]
No distinct pair has valuation at least $K$. Bounding the number of
pairs of valuation exactly $e$ by the number with valuation at least
$e$, and using $|S_Y|\le2Y^s$, proves
\[
 \sum_{N<N'}\varphi^{v_3(N'-N)}
 \le2Y^{2s}\sum_{e\ge1}(\varphi/2)^e\ll Y^{2s}.
\]
\end{proof}

\begin{proposition}[Overlap bound]\label{count:prop:overlap}
For $X\ge1$, $Y\ge3$ and any $G\subset S_Y$,
\[
 \sum_{\substack{N<N'\\N,N'\in G}}|A_N(X)\cap A_{N'}(X)|
 \ll X^gY^{2s+g}.
\]
\end{proposition}
\begin{proof}
For a distinct pair, put $r=N'/N$ and $e=v_3(N'-N)=v_3(r-1)\ge1$.
If $M\in A_N\cap A_{N'}$, then $x=NM\le XY$ and $x,rx\in\D$.
The assignment $M\mapsto x$ is injective. Taking
$n=\lceil\log_3(XY+1)\rceil$ and applying Lemma \ref{count:lem:golden} gives
\[
 |A_N\cap A_{N'}|\le\varphi^{n+e}
 \ll (XY)^g\varphi^e.
\]
Sum this using Lemma \ref{count:lem:valuation}.
\end{proof}

\section{Proof of the power lower bound}

\begin{proof}[Proof of Theorem \ref{thm:power}]
Take $Y=c_0X^\alpha$, where $c_0>0$ is a sufficiently small fixed
constant. We have $0<\alpha<1/6$. To check the latter without numerical
rounding, note that $s<7/11$ since $2^{11}<3^7$, and $g>3/7$ since
$\varphi^7=13\varphi+8>27$. The ratio $(s-g)/(1+g)$ is increasing
in $s$ and decreasing in $g$, so $\alpha<8/55<1/6$.
Consequently $L=\lfloor\log_3X\rfloor\ge3\lceil2\log_3Y\rceil$
for large $X$. Also $X^sY^{s-1}\to\infty$, since $s>1/2$ and
$\alpha<1/6$.

Proposition \ref{count:prop:first} provides a set $G\subset S_Y$ with first
moment at least $X^sY^{2s-1}/64$. By Proposition \ref{count:prop:overlap},
the sum of pairwise overlaps is at most $C_0X^gY^{2s+g}$.
The ratio of this upper bound to the first-moment lower bound is
\[
 64C_0X^{g-s}Y^{1+g}=64C_0c_0^{1+g}.
\]
Choose $c_0$ so that this is at most $1/2$. Bonferroni's inequality gives
\[
 \MC(X)\ge\left|\bigcup_{N\in G}A_N(X)\right|
 \ge\frac1{128}X^sY^{2s-1}
 =\frac{c_0^{2s-1}}{128}X^\beta.
\]

Finally, if $N,MN\in\D$ with $N>0$, choose $B$ with
$3^B>\max(N,MN)$. For every sequence $(\epsilon_j)\in\{0,1\}^{\N_0}$,
\[
 x=N\sum_{j\ge0}\epsilon_j3^{Bj}
\]
and $Mx$ consist of nonoverlapping copies of the finite digit blocks
for $N$ and $MN$. Thus both belong to $\Sigma$. This embeds a scaled
block Cantor set, of dimension $\log2/(B\log3)>0$, into $C(1,M)$.
Hence $\MH(X)\ge\MC(X)$, proving the second bound.
\end{proof}

\section{A delayed digit equation}
Fix $k\ge0$ and put
\[
             m=4k+3,\qquad M=23\cdot3^m-20.
\]
Suppose, for a contradiction, that $Ma=b$ with $a,b\in\D^+$.
Write their finite ternary expansions as
$a=\sum_{j\ge0}a_j3^j$ and $b=\sum_{j\ge0}b_j3^j$, extending both
digit sequences by zeros. For negative indices set $a_j=0$.
The identity $b+20a=23\cdot3^m a$ gives the delayed digits
$x_j=a_{j-m}$ and the recurrence
\begin{equation}\label{family:eq:forward-carry}
 d_0=0,\qquad
 d_{j+1}=\frac{d_j+20a_j+b_j-23x_j}{3}.
\end{equation}

\begin{lemma}\label{family:lem:carry}
Every $d_j$ in \eqref{family:eq:forward-carry} is an integer in
$\I=\{-11,-10,\ldots,10\}$, and $d_j=0$ for all sufficiently
large $j$.
\end{lemma}
\begin{proof}
Set $f_h=20a_h+b_h-23a_{h-m}$. The finite identity
$\sum_{h\ge0}f_h3^h=0$ implies
\[
       d_j=3^{-j}\sum_{h<j}f_h3^h\in\Z,
\]
since the omitted terms are divisible by $3^j$. This formula gives
\eqref{family:eq:forward-carry} and is zero once the partial sum includes
every nonzero term. If $d_j\in\I$, its next numerator lies between
$-34$ and $31$. Since the numerator is divisible by $3$, it lies
between $-33$ and $30$, so $d_{j+1}\in\I$. Induction starts at $0$.
\end{proof}

Encode the last $m$ input digits by the \emph{binary} register
\[
        A_j=\sum_{i=0}^{m-1}a_{j-m+i}2^i.
\]
The oldest digit is in bit $0$, so
\begin{equation}\label{family:eq:register}
 x_j=A_j\bmod2,\qquad
 A_{j+1}=\lfloor A_j/2\rfloor+a_j2^{m-1}.
\end{equation}
This binary encoding is only a convenient label for a queue of
digits. It is not the ordinary ternary value of that queue.
The state space is
\[
             \{0,\ldots,2^m-1\}\times\I.
\]
The initial state is $(0,0)$; a finite witness gives a finite path
back to $(0,0)$ containing an input digit $1$.

\begin{lemma}\label{family:lem:predecessor}
Any positive witness gives a path from $(0,0)$ to $(0,-1)$.
Consequently, no positive witness exists if there is a set of states
containing $(0,-1)$, closed under all predecessors, and omitting
$(0,0)$.
\end{lemma}
\begin{proof}
Because $M\equiv1\pmod3$, the same power of $3$ divides $a$ and $b$.
Dividing it out preserves their allowed digits, so we may assume
$a_0=b_0=1$. The first state after $(0,0)$ is then
$(2^{m-1},7)$, which is different from $(0,0)$.

Consider the first subsequent return to $(0,0)$. A transition into
this state has $a_j=0$ and $A_j=x_j\in\bits$ by
\eqref{family:eq:register}. Equation~\eqref{family:eq:forward-carry} gives
$d_j=23x_j-b_j$. If $x_j=1$, this value is $22$ or $23$, outside
$\I$. If $x_j=0$, the predecessor is $(0,0)$ or $(0,-1)$.
The first possibility is excluded at the first return. Thus the
path passes through $(0,-1)$. A predecessor-closed set containing
that state must contain every earlier state on the path, including
$(0,0)$, giving the asserted contradiction.
\end{proof}

For use in an explicit certificate, all predecessors can be written
without a search. For a current state $(A,d)$ put
$\alpha=\lfloor A/2^{m-1}\rfloor\in\bits$. For each $x,\beta\in\bits$
the potential predecessor is
\begin{equation}\label{family:eq:inverse}
 A^-=2A-\alpha2^m+x,\qquad
 d^-=3d-20\alpha-\beta+23x.
\end{equation}
It is a valid predecessor exactly when $d^-\in\I$.
Here $\alpha$ is the input digit of the forward transition and
$\beta$ is its output digit. Formula~\eqref{family:eq:inverse} enumerates
all choices; automatically $0\le A^-<2^m$.

\section{The predecessor certificate for all \texorpdfstring{$k\ge1$}{k>=1}}
Assume now $k\ge1$. Define
\begin{equation}\label{family:eq:T}
 T_j=\frac{16^j-1}{5}\quad(j\ge1),\qquad
 V=\frac{3T_{k+1}-1}{2}.
\end{equation}
The numbers $T_j$ are integers satisfying $T_{j+1}=16T_j+3$;
their binary expansions are repetitions of the block $0011$ after
padding to length $4j$.

The following finite sets specify the allowed carries in the
certificate:
\begin{align*}
 Z&=\{-10,-9,-4,-3,-1\},&
 O&=\{-8,-7,-5,-4,10\},\\
 P&=\{-2,-1,1,2,7,8,10\},&
 P^+&=P\cup\{4,5\},\\
 Q&=\{-7,-6,-4,-3,2,3,5,6\},&
 R&=\{-10,-9,5,6,8,9\},\\
 S&=\{1,2,4,5,10\},&
 U&=\{-8,-7,-5,-4\},\\
 U^+&=U\cup\{-6,-3\},&
 W&=\{2,3,5,6\},\\
 W^+&=W\cup\{4,7\},&
 H&=\{5,6,8,9\},\\
 I_0&=\{-10,-9\},&
 J_0&=\{-11,-6,-5,-3,-2\},\\
 K_0&=\{-10,-9,-7,-6,8,9\},&
 N&=U.
\end{align*}
The carry-set symbols in this display are local to the certificate.
Distinct row names identify different registers, even when their
carry sets agree.

\begin{table}[htbp]
\centering\small
\caption{Definition of $\Sc_m$. Empty index ranges are omitted.}
\label{family:tab:states}
\begin{tabular}{llll}
\toprule
Row & Index range & Register $A$ & Carry set\\
\midrule
$Z_0$ & & $0$ & $Z$\\
$O_0$ & & $1$ & $O$\\
$P_j$ & $1\le j\le k$ & $T_j$ & $P$\\
$P_{k+1}$ & & $T_{k+1}$ & $P^+$\\
$Q_j$ & $1\le j\le k+1$ & $2T_j$ & $Q$\\
$R_j$ & $1\le j\le k$ & $4T_j$ & $R$\\
$S_j$ & $1\le j\le k$ & $4T_j+1$ & $S$\\
$U_j$ & $1\le j<k$ & $8T_j+1$ & $U$\\
$U_k$ & & $8T_k+1$ & $U^+$\\
$W_j$ & $1\le j<k$ & $8T_j+2$ & $W$\\
$W_k$ & & $8T_k+2$ & $W^+$\\
$H_j$ & $1\le j\le k$ & $T_{j+1}+1$ & $H$\\
$I_0^{\rm row}$ & & $T_{k+1}-1$ & $I_0$\\
$I_1^{\rm row}$ & & $V-1$ & $I_0$\\
$J^{\rm row}$ & & $V$ & $J_0$\\
$K^{\rm row}$ & & $V+1$ & $K_0$\\
$N^{\rm row}$ & & $2T_{k+1}-1$ & $N$\\
\bottomrule
\end{tabular}
\end{table}

\begin{proposition}\label{family:prop:certificate}
The set $\Sc_m$ defined in Table~\ref{family:tab:states} is closed under all
predecessors in \eqref{family:eq:inverse}. It contains $(0,-1)$ and omits
$(0,0)$. Its cardinality is $38k+50$.
\end{proposition}
\begin{proof}
We first determine the most significant bit of every register.
Put $X=16^k\ge16$, so $2^{m-1}=4X$ and $2^m=8X$.
For $j\le k$, $T_j<X/5$. All rows with most significant bit zero
are bounded by
\[
 T_{k+1}+1=\frac{16X+4}{5}<4X;
 \qquad 8T_j+2\le\frac{8X+2}{5}<4X.
\]
The rows $Q_{k+1},I_1^{\rm row},J^{\rm row},K^{\rm row},N^{\rm row}$
have respective registers
\[
 \frac{32X-2}{5},\quad \frac{24X-9}{5},\quad
 \frac{24X-4}{5},\quad \frac{24X+1}{5},\quad \frac{32X-7}{5}.
\]
Each lies in $[4X,8X)$, so its most significant bit is one. This
also verifies that every row defines a state in the prescribed space.

Table~\ref{family:tab:inverse} gives the complete predecessor calculation.
An unannotated destination means its whole carry set. A qualifier
such as ``base $U$'' means the displayed smaller carry set, even
when the destination row uses $U^+$. A dash denotes an empty image.
Every entry allows both values $\beta=0,1$ and retains only
$d^-\in\I$.

\begin{table}[htbp]
\centering\small
\caption{All predecessor images from \eqref{family:eq:inverse}.}
\label{family:tab:inverse}
\begin{tabular}{llll}
\toprule
Current row & $\alpha$ & $x=0$ & $x=1$\\
\midrule
$Z_0$ & $0$ & $Z_0:\{-10,-9,-4,-3\}$ & $O_0$\\
$O_0$ & $0$ & --- & $P_1$ (base $P$)\\
$P_j$, all $j$ & $0$ & $Q_j$ & ---\\
$Q_j$, $j\le k$ & $0$ & $R_j$ & $S_j$\\
$Q_{k+1}$ & $1$ & $J^{\rm row}$ & $K^{\rm row}$\\
$R_j$ & $0$ & --- & $U_j$ (base $U$)\\
$S_j$ & $0$ & $W_j$ (base $W$) & ---\\
$U_j$, $j<k$ & $0$ & --- & $P_{j+1}$ (base $P$)\\
$U_k$ & $0$ & $I_0^{\rm row}$ & $P_{k+1}$ ($P^+$)\\
$W_j$, all $j$ & $0$ & $H_j$ & ---\\
$H_j$ & $0$ & --- & ---\\
$I_0^{\rm row}$ & $0$ & --- & $N^{\rm row}$\\
$I_1^{\rm row}$ & $1$ & --- & ---\\
$J^{\rm row}$ & $1$ & --- & $U_k:\{-7,-6,-4,-3\}$\\
$K^{\rm row}$ & $1$ & $W_k:\{3,4,6,7\}$ & ---\\
$N^{\rm row}$ & $1$ & --- & $I_1^{\rm row}$\\
\bottomrule
\end{tabular}
\end{table}

For the carry column, each entry is the finite calculation
\[
 \{3d-20\alpha-\beta+23x:
               d\text{ in the current carry set},\ \beta\in\bits\}
                    \cap\I.
\]
The register column follows from $T_{j+1}=16T_j+3$ and the two
boundary identities
\begin{equation}\label{family:eq:wrap}
        4T_{k+1}-2^m=V,\qquad 2V-2^m=8T_k.
\end{equation}
For example, the row $Q_{k+1}$ has register $2T_{k+1}$ and
$\alpha=1$. Its predecessor registers are $V,V+1$, and its carry
images are $J_0,K_0$, respectively. The row $J^{\rm row}$ has
$\alpha=1$; for $x=1$ it maps to register $8T_k+1$ and carries
$\{-7,-6,-4,-3\}\subseteq U^+$. These boundary images explain the
extra carries in $P^+,U^+,W^+$.

Thus every possible predecessor lies in a listed row. This verifies
closure for all $k\ge1$ using fixed finite carry sets and the exact
identities \eqref{family:eq:T} and \eqref{family:eq:wrap}. No upper bound on $k$
is used. Since $Z$ contains $-1$ and excludes $0$, and no other row
has register zero, the two required inclusions follow.

For the cardinality, the registers are distinct. In the interior,
\[
 T_j<2T_j<4T_j<4T_j+1<8T_j+1<8T_j+2<T_{j+1},
\]
and $T_{j+1}+1$ lies strictly between $T_{j+1}$ and $2T_{j+1}$.
At the final boundary, writing $T=T_{k+1}\ge51$, the additional
order is
\[
       T-1<T<T+1<V-1<V<V+1<2T-1<2T.
\]
The register $1$ lies below $T_1=3$. Summing the sizes of the carry
sets in Table~\ref{family:tab:states} now gives $38k+50$.
\end{proof}

Lemma~\ref{family:lem:predecessor} and Proposition~\ref{family:prop:certificate}
exclude a positive witness for every $k\ge1$.

\section{The base case and the sporadic exceptional set}
For $k=0$ we have $m=3$ and $M=601$. The small certificate in
Table~\ref{family:tab:base} replaces Table~\ref{family:tab:states}.

\begin{table}[htbp]
\centering\small
\caption{A predecessor-closed certificate of $65$ states for $m=3$.}
\label{family:tab:base}
\begin{tabular}{cl}
\toprule
Register $A$ & Carry values\\
\midrule
$0$ & $\{-10,-9,-4,-3,-1,3,4,6,7,8,9\}$\\
$1$ & $\{-8,-7,-6,-5,-4,-3,10\}$\\
$2$ & $\{-10,-9,3,4,6,7\}$\\
$3$ & $\{-10,-9,-2,-1,1,2,4,5,7,8,10\}$\\
$4$ & $\{-11,-6,-5,-3,-2,8,9\}$\\
$5$ & $\{-10,-9,-8,-7,-6,-5,-4,8,9\}$\\
$6$ & $\{-7,-6,-4,-3,2,3,5,6\}$\\
$7$ & $\{-10,-9,-8,-7,-5,-4\}$\\
\bottomrule
\end{tabular}
\end{table}

Indeed, for each entry $(A,d)$ and $x,\beta\in\bits$, put
\[
 \alpha=\lfloor A/4\rfloor,\qquad
 A^-=2A-8\alpha+x,\qquad d^-=3d-20\alpha-\beta+23x.
\]
Whenever $-11\le d^-\le10$, the value $d^-$ occurs in the row
$A^-$ of Table~\ref{family:tab:base}. This is a direct check of $260$
integer expressions. The table contains $(0,-1)$ and omits $(0,0)$,
so Lemma~\ref{family:lem:predecessor} also excludes a witness for $k=0$.

\begin{proof}[Proof of Theorem~\ref{thm:family}]
The certificates prove $M\notin\mathcal V$ for every $k\ge0$.
Since $M\equiv1\pmod3$, it does not belong to $\Obloc$: an element of
$\Obloc$ coprime to $3$ would have to be congruent to $2$ modulo $3$.
For $m\ge3$,
\[
       2\cdot3^{m+2}<23\cdot3^m-20<\tfrac32\,3^{m+3}.
\]
These are the upper and lower endpoints of two consecutive intervals
in the definition of $\Obreal$, so $M\notin\Obreal$. Hence $M\in\F$.
The numbers $621\cdot3^{4k}-20$ are distinct and unbounded, proving
that $\F$ is infinite.
\end{proof}

\section{An improved exponent without an epsilon loss}\label{ep:section}

This section is new in version 2. The original power bound and the
exceptional-family proof above are retained with their original numbering.
Write $s=\log_3 2$, $\varphi=(1+\sqrt5)/2$, and $g=\log_3\varphi$.

\begin{theorem}\label{ep:main}
There are absolute constants $c_*>0$ and $X_0$ such that
\[
 \MC(X)\ge c_*X^{\beta_*}\quad(X\ge X_0),\qquad
 \beta_*={ (1+s)^2\over7-5s}=0.691726633148796\ldots.
\]
The same lower bound holds for $\MH(X)$. No subtraction of an arbitrary
positive epsilon from $\beta_*$ is required.
\end{theorem}

\subsection{A boundary estimate for the mean Fourier discrepancy}
Use $F_l,f_l,\D_l$ from Section 2. Their second and fourth moments are
$2^{-l}$ and $(3/8)^l$. Let $T_D$ be all nonzero reduced fractions on
$\T$ with denominator at most $D$ coprime to 3, and put
$S_D(L)=\sum_{z\in T_D}f_L(z)$ and $h=\lceil2\log_3D\rceil$.
Multiplication by 3 permutes $T_D$, whose separation is at least $D^{-2}$.
The sampling lemma gives, uniformly for $0\le m\le h$,
\[
 \sum_{T_D}f_m^2\ll D^2 2^{-m},\qquad
 \sum_{T_D}f_m^4\ll D^2(3/8)^m.
\]
Indeed $3^h<3D^2$, so the degrees in both applications are $O(D^2)$.
Set
\[
 a_*={3s-1\over4},\quad A_*={2(3-s)\over3s-1},\quad
 \eta_0=3-4s,\quad r_*={3-s\over2}=a_*A_*.
\]
Here $4<A_*<6$, $0<\eta_0<s$, and $r_*>1$.
For $2h\le L<3h$ write $l=L-2h$ and apply Holder's inequality
with exponents $2,4,4$ to blocks of lengths $h,h,l$. This yields
\begin{equation}\label{ep:farey}
 S_D(L)\ll D^2 2^{-h/2}(3/8)^{(h+l)/4}
 \le C D^{(3+s)/2}3^{-a_*L}.
\end{equation}
The last inequality uses $s-1<0$ and $h\ge2\log_3D$.
For $L\ge3h$, dropping extra factors instead gives
\begin{equation}\label{ep:long}
 S_D(L)\le S_D(3h)\ll D^{\eta_0}.
\end{equation}
For each fixed $D$, also $S_D(L)\to0$: every nonzero fraction has a
finite orbit under multiplication by 3, with all factors
$|(1+\exp(2\pi iz))/2|<1$.

Let $S_Y=\{N\le Y:N\in\D,\ N\equiv1\pmod3\}$ and
$R_N(L)=\sum_{u=1}^{N-1}f_L(u/N)$.
\begin{lemma}\label{ep:mean}
There is an absolute $C_*$ such that, for every fixed $H\ge0$ and
every integer choice $L=L(Y)\ge A_*\log_3Y+H$,
\[
 \limsup_{Y\to\infty}Y^{-s}\sum_{N\in S_Y}R_N(L)
 \le C_*3^{-a_*H}.
\]
\end{lemma}
\begin{proof}
Fixing the high ternary digits shows that
$T_Y(d)=\#\{N\in S_Y:d\mid N\}\le4(Y/d)^s$ for $1\le d\le Y$:
choose a low block with $3^j\le d<3^{j+1}$; its values lie in an
interval of length less than $d$, hence contribute at most one multiple
for each high block. Reducing $u/N$ gives the exact identity
\[
 \sum_{N\in S_Y}R_N(L)=
 \sum_{\substack{2\le d\le Y\\(d,3)=1}}T_Y(d)
       \sum_{\substack{1\le u<d\\(u,d)=1}}f_L(u/d).
\]
In a fixed dyadic range $D/2<d\le D$, its contribution divided by
$Y^s$ is $O(D^{-s}S_D(L))$. One may take $D$ to be powers of 2,
with the final partial range bounded by $D<2Y$.
For ranges with $3h\le L$, equation~\eqref{ep:long} supplies the
summable majorant $CD^{\eta_0-s}$. Each fixed range tends to zero,
so dominated convergence makes their total $o(1)$.
For all remaining ranges, $2h\le L$ once $Y$ is sufficiently large,
because $A_*>4$. Equation~\eqref{ep:farey} gives
\[
 D^{-s}S_D(L)\ll D^{r_*}3^{-a_*L}
 \le 3^{-a_*H}(D/Y)^{r_*}.
\]
The sum of $(D/Y)^{r_*}$ over these dyadic ranges is bounded absolutely.
This proves the stated limsup. In particular the estimate asserts an
arbitrarily small fixed proportion, not necessarily $o(Y^s)$ at fixed $H$.
\end{proof}

\subsection{Affine tails and a fixed residue class}
\begin{lemma}\label{ep:affine}
If $r\in1+3\Z_3$, $r\ne1$, $e=v_3(r-1)$, and $\xi\in\Z_3$, then
\[
 \#\{z\in\D_m:[rz+\xi]_m\in\D_m\}
 \le H_m(e)\le\varphi^m(2/\varphi)^e,
\]
where $[x]_m$ is the residue in $[0,3^m)$ and $H_m(e)$ counts the
independent sets of the graph on $0,\ldots,m-1$ with edges of length $e$.
\end{lemma}
\begin{proof}
Write $r=1+3^eu$, with $u$ a unit. In the addition
$z+3^euz+\xi$, let $B_j$ be the sum of the incoming carry,
the digit of $\xi$, and the shifted digit of $uz$.
On an admissible path the outgoing carry is
$\lfloor(B_j+1)/3\rfloor$, independent of the admissible choice of
$z_j$. Residues $B_j=0,1,2\pmod3$ allow, respectively, two choices,
only $z_j=0$, or only $z_j=1$.
At position $i+e$, after fixing $z_{<i}$, the output digit is
$z_{i+e}+u_0z_i+Q_i(z_{<i})\pmod3$, with $u_0\in\{1,2\}$.
Thus one of the two choices at $i$ forces the digit at $i+e$,
independently of intermediate admissible digits. Backward induction
on the remaining positions bounds the tree by the independent-set
recursion: at a free vertex one either removes that vertex, or removes
it together with its future neighbor; an already forced vertex has
only one branch. Discarding further restrictions only enlarges the
bound. The graph splits into $e$ chains, including empty chains.
A chain of length $l$ has $F_{l+2}\le2\varphi^{l-1}$ independent sets,
also for $l=0$. Multiplying proves the second inequality.
\end{proof}

Fix $c=112=(11011)_3$ and $d=128$ throughout the rest of the proof.
Define
\[
 C_d(c)=\#\{a,b,a',b'\in\D_d:
              c(a-a')+(b-b')\equiv0\pmod{3^d}\}.
\]
For every right-hand side the corresponding count is at most $C_d(c)$,
by Cauchy--Schwarz applied to the autocorrelation of
$h(z)=\#\{(a,b)\in\D_d^2:ca+b\equiv z\pmod{3^d}\}$.
Also $C_d(c)\ge(16/3)^d$.

The following finite certificate specifies the constant used here
without a numerical spectral-radius approximation. On states
$-56,\ldots,56$ put $T_{u,v}=w(3v-u)$, where
\[
 w(0)=4,\quad w(\pm1)=2,\quad w(\pm111)=1,\quad
 w(\pm112)=2,\quad w(\pm113)=1
\]
and all other weights vanish. These are precisely the digit-difference
multiplicities. The state interval is closed under every integral
transition. Starting with $v_0(u)=\mathbf1_{u=0}$, iterate
$v_{j+1}(v)=\sum_u v_j(u)T_{u,v}$. Then
$C=C_d(c)=\sum_u v_{128}(u)$, with unrestricted final carry, equals
\[
\begin{split}
 C={}&2848039325767103655925435712986932980549021547392299357785\\[-2pt]
     &\hspace{8mm}070566911039704321748559183306903592
\end{split}
\]
where the two lines are concatenated decimal digits.
The exact integer checks supplied with the article give
\begin{equation}\label{ep:cert}
 C^{1161}3^{148608}\varphi^{491392}<2^{937216},\qquad
 2^{20805}<3^{13127}.
\end{equation}
These comparisons require integers only: for $m=491392$, write
$\varphi^m=F_m\varphi+F_{m-1}$, multiply by
$C^{1161}3^{148608}$ to obtain $A\varphi+B$, and check
$D=2^{937217}-A-2B>0$ and $D^2>5A^2$.
The executable certificate recomputes the matrix both forwards and
backwards and performs these comparisons exactly.
Set $\lambda=C^{1/d}$ and $\tau=\log_3\lambda$.

\begin{lemma}\label{ep:pairs}
Let $Y=3^k$, $k>d$, and $G\subset\{N\in\D_k:N\equiv1\pmod{3^d}\}$.
For $Y<X$ put
$A_N(X)=\{1\le M\le X:M\equiv c\pmod{3^d},\ NM\in\D\}$.
Then
\[
 \sum_{N<N'\in G}|A_N(X)\cap A_{N'}(X)|\ll_d X^gY^{\tau-g}.
\]
\end{lemma}
\begin{proof}
For four prefixes $N',x,N,y$ of length $m\ge d$ with low blocks
$1,c,1,c$ and $N'x\equiv Ny\pmod{3^m}$, adding four blocks of
length $d$ reduces the next congruence to
\[
 R+c(a-a')+(b-b')\equiv0\pmod{3^d},\qquad R=(N'x-Ny)/3^m.
\]
The quadratic terms vanish since $m\ge d$. There are at most $C$
extensions. If the first distinct denominator digit occurs at $e$,
there are at most $4^{e-d}$ common prefixes and two choices at $e$.
Applying full blocks and bounding a remaining partial block trivially
by $16^d$ proves that the number of length-$k$ quadruples is
$O_d(4^e\lambda^{k-e})$.

Choose $n\ge k$ with $X<3^n\le3X$. For each pair, the map
$M\mapsto[NM]_n$ is injective, since $N$ is a unit. Conditional on a
prefix $x_0$ of length $k$, its remaining digits satisfy an affine
condition with ratio $r=N'/N$ and shift
$\xi=(rx_0-[rx_0]_k)/3^k$. Lemma~\ref{ep:affine} bounds these tails.
Summing over $e$ therefore gives
\[
 O_d\left(\varphi^{n-k}\lambda^k
       \sum_{e\ge d}\left({8\over\varphi\lambda}\right)^e\right).
\]
The ratio is at most $3/(2\varphi)<1$, since $\lambda\ge16/3$.
This proves the assertion, including arbitrarily close distinct pairs.
\end{proof}

\subsection{The first moment and completion of the proof}
Put
\[
 \alpha_*={1\over A_*-1}={3s-1\over7-5s},\qquad
 Y=3^k=3^{\lfloor\log_3(c_0X^{\alpha_*})\rfloor},
\]
where $c_0>0$ will be fixed before $X\to\infty$.
The shell
$\mathcal S_{k,d}=\{N\in\D_k:N\equiv1\pmod{3^d},\ N\ge3^{k-1}\}$
has exactly $2^{-d-1}Y^s$ members, all in $[Y/3,Y/2)$.
Choose a fixed $H$ with $C_*3^{-a_*H}<2^{-d-5}$, and then
$c_0\le3^{-(H+d+4)/(A_*-1)}$. For
$L_* =\lfloor\log_3(XY/3)\rfloor-d$ one has
$L_*-A_*\log_3Y\ge H+2$.
Lemma~\ref{ep:mean} and Markov's inequality show that at least half
the shell, denoted $G$, has $R_N(L_*)\le1/2$ for all sufficiently large $X$.

For $L_N=\lfloor\log_3(NX)\rfloor$, count multiples of $N$ of the form
$z=c+3^dw$, $w\in\D_{L_N-d}$. Orthogonality, permutation of residues
by $3^d$, and monotonicity of $f_l$ give at least
\[
 {2^{L_N-d}\over N}(1-R_N(L_N-d))
 \ge2^{-d-2}X^sN^{s-1}
\]
such multiples. They are positive and less than $NX$, and $z/N$ lies
in the required residue class. Consequently
\begin{equation}\label{ep:first}
 \sum_{N\in G}|A_N(X)|\gg_d X^sY^{2s-1}.
\end{equation}
The exact inequalities~\eqref{ep:cert} imply
$\alpha_*<\alpha_0=1161/5000$ and
$g-s+\alpha_0(1+\tau-g-2s)<0$.
The coefficient $1+\tau-g-2s$ is positive, since
$\tau\ge4s-1$ and $2s>g$. Hence
\[
 g-s+\alpha_*(1+\tau-g-2s)<0.
\]
Lemma~\ref{ep:pairs} is thus $o(1)$ times~\eqref{ep:first}.
Bonferroni's inequality gives
\[
 \MC(X)\ge\left|\bigcup_{N\in G}A_N(X)\right|
 \gg X^sY^{2s-1}\gg X^{s+(2s-1)\alpha_*}=X^{\beta_*}.
\]
The positive-witness inclusion proved above transfers the bound to
$\MH$. This completes the proof of Theorem~\ref{ep:main}.
The constants can be very small and the threshold very large; no
practical search threshold or sharp asymptotic is asserted.

\section{Verification and scope}

For the new bound, \texttt{code/endpoint/run\_checks.py} independently
recomputes the exact 113-state energy forwards and backwards, checks
the integer sign certificates, and tests block lifting, affine tails,
and endpoint algebra. These are finite inputs and consistency checks;
the limiting estimates are proved in Section~\ref{ep:section}.
The underlying archived arguments had internal reviews; the new
English section received a focused integration review. The retained
version-1 proof bodies are not represented as freshly re-reviewed.


All estimates for the power lower bound, including the required
uniformity and limiting arguments, are proved in the text. The
finite-word estimate counts every locally admissible word; it does
not assume an infinite continuation. The accompanying elementary
program \texttt{verify\_power\_gain\_lemmas.py} checks exact finite
digital energies, the finite-word bound for unit ratios modulo $3^n$
through $n=9$, and the stated divisor-cylinder and valuation-bucket
bounds at finite scales. These checks supplement the proofs. They
cannot establish their asymptotic conclusions by themselves.

The nonrepresentation proof consists of the delayed-digit argument
and the explicit predecessor certificates. Its parameter-dependent
part is a symbolic calculation, rather than an extrapolation from
finitely many tested values of $k$. The program
\texttt{verify\_\allowbreak{}parametric\_\allowbreak{}certificate.py}
checks every carry image in
Table~\ref{family:tab:inverse}. Splitting the combined rows into
interior and boundary variants gives $18$ symbolic rows and $384$
choices of carry and binary digits. The register identities are checked
as affine polynomials in formal variables $X=16^k$ and $Y=T_j$, with
exact rational coefficients. The most significant bit inequalities
are proved in Proposition~\ref{family:prop:certificate}. The program
also checks all $260$ predecessor choices in the separate $k=0$
certificate and makes redundant integer instantiations to detect
indexing or transcription errors. Both programs use only the Python
standard library and exact integer or rational arithmetic.

Theorems~\ref{ep:main} and~\ref{thm:family} give a power lower bound
and an explicit infinite sporadic family, respectively. They do not
give a full asymptotic, a density theorem, or a classification of all
nonrepresentable multipliers. The known positive-witness construction
is sufficient, not necessary, for positive Hausdorff dimension.

ChatGPT was used in developing and editing the arguments, implementing
checks, and conducting internal AI-assisted reviews from several
mathematical perspectives. Those reviews are not external peer review
or proof-assistant verification. Authorship is attributed solely to
Timofei Iuzhakov.

\begin{thebibliography}{9}
\bibitem{Lagarias}
J. C. Lagarias,
\emph{Ternary expansions of powers of 2},
J. London Math. Soc. (2) \textbf{79} (2009), 562--588.
\href{https://doi.org/10.1112/jlms/jdn080}{doi:10.1112/jlms/jdn080}.
\bibitem{AL}
W. C. Abram and J. C. Lagarias,
\emph{Intersections of multiplicative translates of 3-adic Cantor sets},
J. Fractal Geom. \textbf{1} (2014), 349--390.
\url{https://ems.press/content/serial-article-files/37215}.
\bibitem{ABL}
W. C. Abram, A. Bolshakov and J. C. Lagarias,
\emph{Intersections of multiplicative translates of 3-adic Cantor sets II:
two infinite families},
Experimental Math. \textbf{26} (2017), 468--489.
\href{https://doi.org/10.1080/10586458.2016.1205532}{doi:10.1080/10586458.2016.1205532}.
\bibitem{BMRS}
J. H. Bai, J. Meleshko, S. Riasat and J. Shallit,
\emph{Quotients of palindromic and antipalindromic numbers},
Integers \textbf{22} (2022), A96.
\href{https://math.colgate.edu/~integers/w96/w96.pdf}{Journal article};
\href{https://arxiv.org/abs/2202.13694}{arXiv:2202.13694}.
\bibitem{ADRS}
K. Anders, M. L. Dawsey, B. Reznick and S. Sisneros-Thiry,
\emph{Representations of integers as quotients of sums of distinct powers
of three}, preprint, 2023.
\url{https://arxiv.org/abs/2308.07252}.
\bibitem{Konyagin}
S. Konyagin,
\emph{Arithmetic properties of integers with missing digits:
distribution in residue classes},
Period. Math. Hungar. \textbf{42} (2001), 145--162.
\href{https://doi.org/10.1023/A:1015256809636}{doi:10.1023/A:1015256809636}.
\bibitem{MS}
N. G. Moshchevitin and I. D. Shkredov,
\emph{On the multiplicative properties modulo $m$ of numbers with missing
digits}, Math. Notes \textbf{81} (2007), 338--355.
\url{https://www.mi-ras.ru/~ishkredov/Papers/9.pdf}.

\bibitem{Jandhyala}
V. Jandhyala,
\emph{Infinitely many integers that are not quotients of zero-free
balanced ternary numbers}, preprint, version~3, 24 October 2026.
\href{https://github.com/jvvk/mathematics/tree/main/balanced-ternary-quotients}{Author's paper and companion repository}.
\end{thebibliography}
\end{document}
